\documentclass{article}
\usepackage{amsmath,amsthm}
\newtheorem{proposition}{Proposition}

\begin{document}

\section{Wages}

Four workers, observed in 2019 and 2020. One 2020 wage is a coding error
in the payroll system; we drop it rather than repair it, so the drop
stays visible.

\begin{table}[t]
  \centering
  \caption{Log wage moments by year}\label{tab:moments}
  \input{tables/moments.tex}

  \footnotesize Notes: non-positive wages are dropped. The variance is
  the population variance of log wages (dividing by $n$, not $n-1$).
\end{table}

\begin{proposition}\label{prop:units}
  The variance of log wages in Table~\ref{tab:moments} does not depend on
  the currency unit: paying every worker $c > 0$ times their wage leaves
  it unchanged.
\end{proposition}

\begin{proof}
  Rescaling adds $\log c$ to every log wage, and a variance is unchanged
  by adding a constant.
\end{proof}

\begin{figure}[h]
  \centering
  \fbox{payroll extract $\to$ positive wages $\to$ moments by year}
  \caption{How the numbers in Table~\ref{tab:moments} are made.}\label{fig:flow}
\end{figure}

\end{document}
